\section{Language Identification}
\label{sec:lid}

Language Identification (LID) is the task of determining the language which is spoken in a given utterance. 
This has several important applications: 
despite much work on multilingual speech recognition~\citep{burget2010multilingual, lin2009study,heigold2013multilingual,bourlard2011current,cho2018multilingual,toshniwal2018multilingual,kannan2019mult,li2019bytes,pratap2020massively}, many deployed systems are still trained on data in a single language despite the need to  transcribe speech in different languages. 
It is therefore crucial to route utterances to the correct system and this routing depends on an LID model.
Moreover, much work on mining speech data from the web relies on LID, including some of the most recent large-scale weakly-supervised work~\citep{radford2022whisper}.
Training a language identification system requires speech data for which the spoken language is known, however, the most diverse public corpora span no more than about 100 languages~\citep{valk2020voxlingua,conneau2022fleurs}.

In this section, we first describe our methodology to build LID models (\textsection\ref{sec:lid-setup}), then evaluate the feasibility of training LID systems using \MMSU{} and \GRN{} data compared data from existing corpora (\textsection\ref{sec:lid-existing-data}) and then build LID models with 40x more languages compared to existing systems (\textsection\ref{sec:lid-scaling}).\footnote{Note that we use \MMSU{} instead of \MMS{} because it supports more languages and we do not require the data to be paired with transcriptions for LID (\textsection\ref{sec:data-paired-source}).}



\subsection{Training Setup}
\label{sec:lid-setup}

We train models by fine-tuning the \modelnameb{1} pre-trained model (\textsection\ref{sec:pretrain-setup}) for language identification.
This is done by stacking a linear classifier on top of the pre-trained model, which maps to the set of possible languages for a particular task, followed by fine-tuning all parameters, including the pre-trained model.

We optimize models with Adam with exponential decay rates $\beta_1=0.9$ and $\beta_2=0.98$ and a tri-state learning rate schedule where the learning rate is warmed up for the first 10\% of updates, held constant for the next 40\% and then linearly decayed for the remainder of training~\citep{kingma2015adam}.
During development, we experiment with different hyper-parameters and perform final model selection based on development set accuracy.
We experiment with different learning rates ($\Enot{1e-5}$, $\Enot{3e-5}$, $\Enot{3e-6}$, $\Enot{5e-6}$, $\Enot{7e-6}$), training updates (10K, 20K, 30K, 40K, 50K) and batch sizes (1.5min, 3min, 6min). 
We train models on 16 GPUs.

To balance the different languages and corpora during training, we first balance the data of every language in the different corpora, using sampling parameter $\beta_L$. This is followed by balancing each language using the resampled data of the first step, using sampling parameter $\beta_D$ with the sampling distribution outlined in~\textsection\ref{sec:pretrain-w2v}.
We experiment with $\beta_L$ and $\beta_D$ settings 0, 0.3, 0.5, 0.7 and 1.

When we train models on multiple datasets, then we use a development set containing up to 30 minutes of data for each language and we sample an equal amount of data from each corpus.
For models supporting 1K, 2K or 4K languages, we reduce the amount of data per language to 15 min, 7min and 3min to enable faster development.



\subsection{Comparison to Existing Datasets}
\label{sec:lid-existing-data}

\insertResultsLIDFLVLCompare


We first assess the effectiveness of training LID models on \MMSU{} and \GRN{} data compared to existing LID training data (FLEURS, VoxLingua-107).
Performance is evaluated both on the FLEURS and VoxLingua-107 benchmarks.
In this setting, FLEURS and VoxLingua-107 models are at an advantage because there is no domain shift compared to systems trained on \MMSU{} and \GRN{} data.

We are particularly interested in the setting where models trained on existing corpora are evaluated out-of-domain, i.e., the performance of a model trained on VoxLingua-107 evaluated on FLEURS or a model trained on FLEURS evaluated on VoxLingua-107 data, and how this compares to models trained on \MMSU{}/\GRN{}.
To enable a controlled comparison, we report LID accuracy on the 72 languages supported by all considered datasets and train only on data of these languages.


The results (\autoref{fig:lid-fl-vl-comp2}) show that \MMSUboth{} enables LID with slightly lower performance compared to systems trained on existing datasets when both are evaluated out-of-domain: 
on FLEURS evaluation data the gap is 2.1\% compared to a VoxLingua-107 model and on VoxLingua-107 evaluation data \MMSUboth{} trails a FLEURS model by 1.6\%.
Combining \MMSU{} and \GRN{} works particularly well as \GRN{} is more varied which improves performance.
Naturally, models trained on in-domain training data (FLEURS or VoxLingua-107) perform best.
We conclude that \MMSU{} and \GRN{} enable good quality LID models while enabling LID for many more languages as we will demonstrate next.



\subsection{Scaling Language Identification to \numlidlang{} Languages}
\label{sec:lid-scaling}

\insertLIDResultsScaling

Next, we scale spoken language identification from about 100 languages to \numlidlang{} languages by combining \MMSU{}, \GRN{}, FLEURS, and VoxLingua-107 data.
Our primary goal is to understand how accuracy is impacted as models support more and more languages.
We start with 126 languages, the union of the languages in both FLEURS and VoxLingua-107, and then add languages from \MMSUboth{}, roughly doubling the number of languages at every experiment.
Languages are sorted by descending speaker count and we add the most spoken languages first.\footnote{We obtain the number of speakers of each language from \url{https://www.ethnologue.com}}

Performance is evaluated on FLEURS and VoxLingua-107 which is in-domain with respect to the models we train.
To get a sense of how the models perform on domains not seen in the training data, we also evaluate on two other datasets which are out-of-domain with respect to the training data domain:
first, BABEL is conversational telephone speech data and we evaluate on 23 African and Asian languages~\citep{gales2014babel}.\footnote{Amharic, Assamese, Bengali, Cantonese, Cebuano, Georgian, Guarani, Haitian Creole, Igbo, Javanese, Kazakh, Lao, Lithuanian, Dholuo, Mongolian, Pashto, Swahili, Tagalog, Tamil, Telugu, Turkish, Vietnamese,  Zulu. 
We evaluate on test utterances longer than 10 seconds.}
Second, VoxPopuli~\citep{wang2021voxpopuli} consists of parliamentary speech in 25 languages from the European parliament. 
We evaluate on 2.5 hours of data for each language sampled from the VoxPopuli unlabeled data portion.

\autoref{tab:lid-scale} shows that \modelname{} models scale very well: 
increasing the number of languages from 126 to \numlidlang{} results in a modest performance drop of just 0.3\% on FLEURS and no drop on VoxLingua-107.
Out-of-domain we observe a drop of 3.6\% on BABEL and 0.2\% on VoxPopuli.
The results are also competitive to models trained only on in-domain data: 
On FLEURS evaluation data, the \numlidlang{} language \modelname{} model performs 1\% better than the baseline trained only on FLEURS data and is only 0.2\% behind the baseline trained on both FLEURS and VoxLingua-107.
On VoxLingua-107 evaluation data, the gap is 0.4-0.8\%.
This shows that scaling LID to \numlidlang{} can result in models with competitive performance to models trained on much fewer languages.
